\documentclass[10pt,a4paper]{amsart}
\usepackage[margin=19mm]{geometry}
\usepackage{amsmath,amssymb,mathtools}
\usepackage[scr=rsfs,cal=euler]{mathalfa}
\usepackage{xfrac}
\usepackage[unicode,hidelinks,bookmarks=false]{hyperref}
\newcommand{\ie}{i.e.\ }
\newcommand{\cf}{cf.\ }
\newcommand{\resp}{respectively\ }
\newcommand{\Subsection}{\subsection}
\providecommand{\nobreakdash}{\nobreak\hbox{-}}
\setlength{\emergencystretch}{2em}
\multlinegap0pt
\usepackage{mathtools}
\usepackage{amssymb}
\usepackage[normalem]{ulem}
\newtheorem{lem}{Lemma}
\title{Inversion of adjunction for quotient singularities III: semi-invariant case}
\author{Yusuke Nakamura, Kohsuke Shibata}
\date{Source: arXiv:2312.05808v2 (CC BY 4.0)}
\begin{document}
\maketitle

\noindent\emph{Source and licence.} Inversion of adjunction for quotient singularities III: semi-invariant case. Version arXiv:2312.05808v2 (CC BY 4.0), distributed by arXiv under the Creative Commons Attribution 4.0 International licence, http://creativecommons.org/licenses/by/4.0/, as recorded in the arXiv licence metadata for this exact version and shown on the abstract page https://arxiv.org/abs/2312.05808v2. Full licence notice: \texttt{licenses/arxiv-2312.05808-CC-BY-4.0.txt}.\par\smallskip

\noindent\textit{Editorial scope.} Counted unit: the article's lem lem:L4 (source0.tex lines 1382--1385). The article's complete original proof of it is delivered with it (source0.tex lines 1386--1472, one \texttt{proof} environment of 87 lines). No mathematics is written, repaired or re-derived: the statement and the proof are the article's own lines, in the article's own notation, and no retained line is substituted or re-flowed.  Retained uncounted as the source-local context the proof needs: lines 827--842; lines 1126--1142.  Nothing was omitted from any retained span, so the package carries no omission record.  No retained line contains a citation, so the wrapper carries no bibliography and no bibliography entry.  No retained line contains a figure environment, an image-inclusion command or drawing code, so no image is delivered and the package declares no asset dependency.  Editorial formatting only, in addition: an AMS \texttt{amsart} preamble that restates the article's own declarations and macros used by the retained lines, namely the environments \texttt{lem} on separate counters, together with the standard packages those lines need.  The retained environments print in this document as Lemma 1 in order of appearance, whereas the article prints the same items with the numbering of its own full document.  The complete native source of the article as distributed in the arXiv e-print for this exact version is bundled unchanged as \texttt{sources/151255/source0.tex}.
\begin{lem}\label{lem:lambda}
The following assertions hold. 
\begin{enumerate}
\item 
For any $\delta \in C_{\gamma}$, the composition $R \xrightarrow{\delta} R \xrightarrow{\lambda ^* _{\gamma}} R[t^{1/d}]$ coincides with the composition $R \xrightarrow{\lambda ^* _{\gamma}} R[t^{1/d}] \xrightarrow{\delta} R[t^{1/d}]$. 
Here, $R[t^{1/d}] \xrightarrow{\delta} R[t^{1/d}]$ is the $k[t^{1/d}]$-ring homomorphism induced by $R \xrightarrow{\delta} R$.

\item
We have 
$\lambda^* _{\gamma} \left( R^G \right) \subset R^{C_{\gamma}}[t]$. 

\item 
If $f \in R$ is a $G$-semi-invariant element, then 
for any $\delta \in G$, we have $\lambda ^* _{\delta}(f)t^{-w_{\delta}(f)} \in R[t]$.
\end{enumerate}
\end{lem}

\begin{lem}\label{lem:L}
\begin{enumerate}
\item
There is a natural morphism $\Omega ^n _{\overline{B}^{(\gamma)}_{1/d}/k[t^{1/d}]} \to L_{\overline{B}^{(\gamma)}_{1/d}}$, and we have
\[
\operatorname{Im} \left ( \Omega ^n _{\overline{B}^{(\gamma)}_{1/d}/k[t^{1/d}]} \to L_{\overline{B}^{(\gamma)}_{1/d}} \right)  = \operatorname{Fitt}^n \Bigl ( \Omega _{\overline{B}^{(\gamma)}_{1/d}/k[t^{1/d}]} \Bigr ) \otimes _{\mathcal{O}_{\overline{B}^{(\gamma)}_{1/d}}} L_{\overline{B}^{(\gamma)}_{1/d}}.
\]  

\item
There is a natural morphism $t^{r \cdot w_{\gamma}({\bf f})} \Bigl ( \omega^{[r]} _{B_{k[t^{1/d}]}/k[t^{1/d}]} \big| _{\overline{B}^{(\gamma)}_{1/d}} \Bigr ) \to L_{\overline{B}^{(\gamma)}_{1/d}}^{\otimes r}$, 
and we have
\[
\operatorname{Im} \left (t^{r \cdot w_{\gamma}({\bf f})} \Bigl ( \omega^{[r]} _{B_{k[t^{1/d}]}/k[t^{1/d}]} \big| _{\overline{B}^{(\gamma)}_{1/d}} \Bigr ) \to L_{\overline{B}^{(\gamma)}_{1/d}}^{\otimes r} \right)
= t^{r \cdot \operatorname{age}(\gamma)} L_{\overline{B}^{(\gamma)}_{1/d}}^{\otimes r}. 
\]
\end{enumerate}
\end{lem}

\begin{lem}\label{lem:L4}
There exists a $C_{\gamma}$-invariant ideal sheaf $\mathfrak{n}'_{p}$ on $\overline{B}^{(\gamma)}$ such that 
$\mathfrak{n}''_{\overline{B}^{(\gamma)}_{1/d}} = \mathfrak{n}'_{p} \mathcal{O}_{\overline{B}^{(\gamma)}_{1/d}}$. 
\end{lem}

\begin{proof}
Take elements $g_1, \ldots , g_{\ell} \in R^{C_{\gamma}}$ satisfying $R^{C_{\gamma}} = k[g_1, \ldots, g_{\ell}]$. 
For a subset $J \subset \{ 1, \ldots, \ell \}$ with $\# J = N-c$, we define $\Delta _J \in R[t]$ as the determinant of the Jacobian matrix (of size $N$) with respect to $\lambda ^* _{\gamma} (f_i) t^{- w_{\gamma}(f_i)}$ for $1 \le i \le c$ and $g_i$ for $i \in J$, and $\partial x_j$ for $1 \le j \le N$. 
Note that $\Delta _J$ is defined up to sign.
We define an ideal $\mathfrak{n}'_{p} \subset R[t] / I_{\overline{B}^{(\gamma)}}$ as 
\[
\mathfrak{n}'_{p} := \bigl ( \overline{ \Delta _J} \ \big| \ \text{$J \subset \{ 1, \ldots, \ell \}$ with $\# J = N-c$} \bigr ), 
\]
where $\overline{ \Delta _J}$ denoted the image of $\Delta _J \in R[t]$ in $R[t] / I_{\overline{B}^{(\gamma)}}$. 
For any $\delta \in C_{\gamma}$, by Lemma \ref{lem:lambda}(1), we have 
\[
\delta \left( \lambda ^* _{\gamma} (f_i) t^{- w_{\gamma}(f_i)} \right) 
= 
\xi^{d w_{\delta} (f_i)} \lambda ^* _{\gamma} (f_i) t^{- w_{\gamma}(f_i)}. 
\]
Furthermore, note that $\delta$ maps $x_i$'s to their $k$-linear independent $k$-linear combinations. 
Therefore, for each $J$, we have 
\begin{align*}
&
\xi^{d w_{\delta} ({\bf f})} \cdot \Delta _J \cdot dx_1 \wedge \cdots \wedge dx_N \\
&=
\xi^{d w_{\delta} ({\bf f})} \cdot 
\bigl( \wedge_{i \in J} \ d g_i \bigr) 
\wedge 
\left( \wedge _{1 \le i \le c}\ d \left( \lambda ^* _{\gamma} (f_i) t^{- w_{\gamma}(f_i)} \right) \right) \\
&=
\bigl( \wedge_{i \in J}\ d (\delta (g_i)) \bigr) 
\wedge 
\left( \wedge _{1 \le i \le c} \ d \left( \delta \left( \lambda ^* _{\gamma} (f_i) t^{- w_{\gamma}(f_i)} \right) \right) \right) \\
&=
\delta \left( 
\bigl( \wedge_{i \in J}\ d g_i \bigr) 
\wedge 
\left( \wedge _{1 \le i \le c} \ d \left( \lambda ^* _{\gamma} (f_i) t^{- w_{\gamma}(f_i)} \right) \right)
\right)\\
&=
\delta \left(
\Delta _J \cdot dx_1 \wedge \cdots \wedge dx_N
\right) \\
&=
\delta \left( \Delta _J \right) 
\cdot d (\delta (x_1)) \wedge \cdots \wedge d (\delta (x_N)) \\
&= 
\delta \left( \Delta _J \right) \cdot s \cdot dx_1 \wedge \cdots \wedge dx_N 
\end{align*}
for some $s \in k^{\times}$. 
Hence, we have $\delta (\Delta _J) = s' \cdot \Delta _J$ for some $s' \in k$. 
Hence, we conclude that $\mathfrak{n}'_{p}$ is $C_{\gamma}$-invariant. 

We shall prove that this $\mathfrak{n}'_{p}$ satisfies the assertion $\mathfrak{n}''_{\overline{B}^{(\gamma)}_{1/d}} = \mathfrak{n}'_{p} \mathcal{O}_{\overline{B}^{(\gamma)}_{1/d}}$.  
We set 
\[
S_{\widetilde{B}^{(\gamma)}_{1/d}} := R^{C_{\gamma}}[t^{1/d}] \big/ \left(  I_{\overline{B}^{(\gamma)}} \cap R^{C_{\gamma}}[t] \right) R^{C_{\gamma}}[t^{1/d}]. 
\]
We also use the notation in the proof of Lemma \ref{lem:L}. 
By the proof of Lemma \ref{lem:L}(1), the considered map $\Omega^n _{\widetilde{B}^{(\gamma)}_{1/d}/k[t^{1/d}]} \big|_{\overline{B}^{(\gamma)}_{1/d}} \to L_{\overline{B}^{(\gamma)}_{1/d}}$ is corresponding to 
\begin{align*}
\tag{i}
\left ( \left( \bigwedge ^{c} \left( I_{\overline{B}_{k[t^{1/d}]}}/ I^2_{\overline{B}_{k[t^{1/d}]}} \right) \otimes _{S_{\overline{B}_{k[t^{1/d}]}}} S_{\widetilde{B}^{(\gamma)}_{1/d}} \right)  \otimes _{S_{\widetilde{B}^{(\gamma)}_{1/d}}} \Omega _{S_{\widetilde{B}^{(\gamma)}_{1/d}}/k[t^{1/d}]} ^{N-c} \right) &
\otimes _{S_{\widetilde{B}^{(\gamma)}_{1/d}}} S_{\overline{B}^{(\gamma)}_{1/d}} \\
\longrightarrow \ \ \Omega _{R[t^{1/d}]/k[t^{1/d}]} ^N \otimes _{R[t^{1/d}]} S_{\overline{B}^{(\gamma)}_{1/d}}. & 
\end{align*}

The
$S_{\widetilde{B}^{(\gamma)}_{1/d}}$-module 
\[
\left( \bigwedge ^{c} \left( I_{\overline{B}_{k[t^{1/d}]}}/ I^2_{\overline{B}_{k[t^{1/d}]}} \right) \otimes _{S_{\overline{B}_{k[t^{1/d}]}}} S_{\widetilde{B}^{(\gamma)}_{1/d}} \right) \otimes _{S_{\widetilde{B}^{(\gamma)}_{1/d}}} \Omega _{S_{\widetilde{B}^{(\gamma)}_{1/d}}/k[t^{1/d}]} ^{N-c}
\] is generated by the elements of the form 
\[
\left( \overline{f_1} \wedge \cdots \wedge \overline{f_c} \right) \otimes \left( d \overline{g}_{i_1} \wedge \cdots \wedge d \overline{g}_{i_{N-c}} \right).  
\]
Its image by the map (i) is
\begin{align*}
&d \left( \lambda _{\gamma} ^* (f_1) \right) \wedge \cdots \wedge d \left( \lambda _{\gamma} ^* (f_c) \right) \wedge  dg_{i _1} \wedge \cdots \wedge dg_{i_{N-c}} \\
&=t^{w_{\gamma}({\bf f})} \cdot d \left( \lambda _{\gamma} ^* (f_1) t^{- w_{\gamma} (f_1)} \right) \wedge \cdots \wedge d \left( \lambda _{\gamma} ^* (f_c) t^{- w_{\gamma} (f_c)} \right) \wedge  d g_{i _1} \wedge \cdots \wedge d g_{i_{N-c}}. 
\end{align*}
This is equal to 
\[
t^{w_{\gamma}({\bf f})} \cdot \Delta _J \cdot dx_1 \wedge \cdots \wedge dx_N
\]
for $J = \{ i_1, \ldots , i_{N-c} \}$. 
Therefore, we conclude that the image of the map (i) is equal to 
\[
t^{w_\gamma({\bf f})} \mathfrak{n}'_{p} \left( \Omega _{R[t^{1/d}]/k[t^{1/d}]} ^N \otimes _{R[t^{1/d}]} S_{\overline{B}^{(\gamma)}_{1/d}} \right), 
\]
which completes the proof. 
\end{proof}

\end{document}
